% Capítulo 4 - Metodologia

\chapter{METODOLOGIA}
\label{cap:metodologia}

Este capítulo descreve a metodologia adotada para a montagem, configuração e teste da solução proposta, detalhando os procedimentos, ferramentas e cenários de avaliação.

\section{MONTAGEM DA INFRAESTRUTURA}

A montagem física compreende os seguintes componentes:

\begin{alineas}
    \item Radxa ROCK 4B+ como placa principal;
    \item Fonte de alimentação USB-C 5V/3A compatível;
    \item Cabo de rede \textit{Gigabit Ethernet};
    \item Três SSDs SATA de 2,5 polegadas (480 GB cada);
    \item Três adaptadores USB 3.0 para SATA;
    \item \textit{Hub} USB 3.0 com alimentação externa;
    \item Cartão microSD (32 GB, Classe 10) para sistema operacional;
    \item Dissipador passivo e ventilação para refrigeração.
\end{alineas}

A ordem de montagem segue o fluxo: conexão dos adaptadores USB-SATA ao \textit{hub} USB, conexão do \textit{hub} à ROCK 4B+, inserção dos SSDs nos adaptadores, conexão do cabo de rede e da fonte de alimentação.

\section{INSTALAÇÃO E CONFIGURAÇÃO DO SOFTWARE}

O processo de instalação segue as etapas descritas no Quadro~\ref{quad:instalacao}.

\begin{quadro}[htb]
    \centering
    \legenda[quad:instalacao]{ETAPAS DE INSTALAÇÃO E CONFIGURAÇÃO}
    \small
    \begin{tabularx}{\linewidth}{|c|>{\raggedright\arraybackslash}X|}
        \hline
        \textbf{Etapa} & \textbf{Descrição} \\
        \hline
        \hline
        1 & Gravação da imagem Armbian minimal no cartão microSD com \textit{Raspberry Pi Imager} \\
        \hline
        2 & Inicialização da ROCK 4B+ e configuração de rede (IP fixo) \\
        \hline
        3 & Atualização do sistema com \texttt{apt update \&\& apt upgrade} \\
        \hline
        4 & Instalação do \textit{OpenMediaVault} via script oficial \\
        \hline
        5 & Configuração dos discos: formatação ext4, montagem, pastas compartilhadas \\
        \hline
        6 & Instalação do \textit{Docker} e \textit{Docker Compose} \\
        \hline
        7 & Implantação do \textit{Nextcloud} via \textit{Compose} (PHP, MariaDB, Redis) \\
        \hline
        8 & Configuração do \textit{Tailscale} via \textit{Compose} ou pacote oficial \\
        \hline
        9 & Criação de usuários e permissões no OMV e \textit{Nextcloud} \\
        \hline
        10 & Configuração de compartilhamento SMB/CIFS para rede local \\
        \hline
    \end{tabularx}
    \vspace{0.2cm}
    \fonteautoria
\end{quadro}

\section{PROCEDIMENTOS DE TESTE}

Os testes serão conduzidos em dois cenários: rede local (\textit{Gigabit Ethernet} cabeada) e acesso remoto (via \textit{Tailscale}). Cada teste será executado três vezes para garantir reprodutibilidade, e o valor médio será registrado.

\subsection{Teste de \textit{Throughput} Local}

O \textit{throughput} de leitura e escrita será medido com as seguintes ferramentas e comandos:

\begin{itemize}
    \item \textbf{dd}: Escrita sequencial (\texttt{dd if=/dev/zero of=/mnt/teste/test.img bs=1M count=1024}) e leitura (\texttt{dd if=/mnt/teste/test.img of=/dev/null bs=1M});
    \item \textbf{iperf3}: Medição de \textit{throughput} de rede (\texttt{iperf3 -c <ip> -t 30});
    \item \textbf{rsync}: Transferência de arquivos grandes (1 GB) entre máquina cliente e NAS;
    \item \textbf{hdparm}: Leitura de \textit{cache} e \textit{buffer} do disco (\texttt{hdparm -Tt /dev/sda}).
\end{itemize}

\subsection{Teste de \textit{Throughput} Remoto}

O acesso remoto será testado através do \textit{Tailscale}, com os mesmos procedimentos do teste local, porém utilizando o endereço IP atribuído pela malha \textit{VPN}. A latência será medida com \texttt{ping} e \texttt{mtr}.

\subsection{Teste de Estabilidade}

A estabilidade será avaliada submetendo o sistema a carga contínua com \texttt{stress-ng} por 24 horas, monitorando temperatura (\texttt{sensors}), uso de CPU e RAM (\texttt{htop}), e verificação de quedas de serviço.

\subsection{Monitoramento Térmico}

A temperatura da SBC será registrada a cada 5 minutos durante os testes de estresse, utilizando \texttt{watch -n 300 sensors >> termal.log}. O limite de segurança adotado é 80\,\textdegree C.

\section{CENÁRIOS DE TESTE}

O Quadro~\ref{quad:cenarios} resume os cenários de teste previstos.

\begin{quadro}[htb]
    \centering
    \legenda[quad:cenarios]{CENÁRIOS DE TESTE}
    \small
    \begin{tabularx}{\linewidth}{|c|>{\raggedright\arraybackslash}X|c|c|}
        \hline
        \textbf{Cenário} & \textbf{Descrição} & \textbf{Acesso} & \textbf{Protocolo} \\
        \hline
        \hline
        C1 & Leitura sequencial local (1 SSD) & Local & SMB \\
        \hline
        C2 & Escrita sequencial local (1 SSD) & Local & SMB \\
        \hline
        C3 & Leitura simultânea (3 SSDs) & Local & SMB \\
        \hline
        C4 & Escrita simultânea (3 SSDs) & Local & SMB \\
        \hline
        C5 & \textit{Upload Nextcloud} (100 MB) & Local & HTTPS \\
        \hline
        C6 & \textit{Download Nextcloud} (100 MB) & Local & HTTPS \\
        \hline
        C7 & Leitura sequencial remota & \textit{Tailscale} & SMB \\
        \hline
        C8 & Escrita sequencial remota & \textit{Tailscale} & SMB \\
        \hline
        C9 & \textit{Upload Nextcloud} remoto & \textit{Tailscale} + HTTPS & HTTPS \\
        \hline
        C10 & Estresse 24h (CPU + disco + rede) & Local & Misto \\
        \hline
    \end{tabularx}
    \vspace{0.2cm}
    \fonteautoria
\end{quadro}

\section{FERRAMENTAS UTILIZADAS}

O Quadro~\ref{quad:ferramentas} lista as ferramentas que serão empregadas nos testes e suas respectivas finalidades.

\begin{quadro}[htb]
    \centering
    \legenda[quad:ferramentas]{FERRAMENTAS DE TESTE E MONITORAMENTO}
    \small
    \begin{tabularx}{\linewidth}{|>{\raggedright\arraybackslash}X|>{\raggedright\arraybackslash}X|}
        \hline
        \textbf{Ferramenta} & \textbf{Finalidade} \\
        \hline
        \hline
        \texttt{dd} & Medição de \textit{throughput} de leitura/escrita sequencial em disco \\
        \hline
        \texttt{hdparm} & Medição de velocidade de \textit{cache} e \textit{buffer} do disco \\
        \hline
        \texttt{iperf3} & Medição de \textit{throughput} de rede TCP/UDP \\
        \hline
        \texttt{rsync} & Transferência de arquivos com medição de taxa \\
        \hline
        \texttt{stress-ng} & Geração de carga computacional para testes de estresse \\
        \hline
        \texttt{sensors} & Monitoramento de temperatura da SBC \\
        \hline
        \texttt{htop} & Monitoramento de CPU, RAM e processos em tempo real \\
        \hline
        \texttt{ping} / \texttt{mtr} & Medição de latência e perda de pacotes \\
        \hline
        Medidor de energia & Consumo elétrico (idle e sob carga) \\
        \hline
    \end{tabularx}
    \vspace{0.2cm}
    \fonteautoria
\end{quadro}
